% Preambulo partilhado para todos os resumos de disciplinas.
% Fontes/cores/caixas ficam todas aqui para manter os documentos visualmente consistentes.
% NOTA: não carregar babel/fontenc/inputenc aqui — o template pandoc (via
% -V lang=pt-PT, engine tectonic/xetex) já os configura; duplicá-los causa
% "Option clash for package babel".
\usepackage{amsthm}
\usepackage{graphicx}
\usepackage{tikz}
\usetikzlibrary{automata,positioning,arrows.meta,shapes,trees}
\usepackage[most]{tcolorbox}
\usepackage{xcolor}
\usepackage{fancyhdr}
\usepackage{listings}
\usepackage{enumitem}
\usepackage{booktabs}
\usepackage{longtable}
\usepackage[colorlinks=true,linkcolor=DefBlue,urlcolor=DefBlue,citecolor=DefBlue]{hyperref}
\usepackage{geometry}
\geometry{margin=2.5cm}

\definecolor{DefBlue}{HTML}{1B4F72}
\definecolor{DefBlueBg}{HTML}{EBF2FA}
\definecolor{ExGreen}{HTML}{1E7A46}
\definecolor{ExGreenBg}{HTML}{E9F7EF}
\definecolor{WarnOrange}{HTML}{B9770E}
\definecolor{WarnOrangeBg}{HTML}{FDF2E3}
\definecolor{ExamRed}{HTML}{922B21}
\definecolor{ExamRedBg}{HTML}{FDEDEC}

% Caixas NÃO quebráveis de propósito (nobreak é o default do tcolorbox
% quando "breakable" não é passado): uma caixa nunca deve partir-se entre
% duas páginas — se não couber no espaço que resta, o LaTeX empurra-a
% inteira para a página seguinte. Se um exemplo for demasiado longo para
% caber numa página, divide-o em duas caixas mais pequenas em vez de
% ativar "breakable" outra vez.
% Caixa de Definicao
\newtcolorbox{definicao}[1][]{colback=DefBlueBg,colframe=DefBlue,fonttitle=\bfseries,title={Definição\ #1},left=6pt,right=6pt,top=3pt,bottom=3pt}
% Caixa de Exemplo resolvido
\newtcolorbox{exemplo}[1][]{colback=ExGreenBg,colframe=ExGreen,fonttitle=\bfseries,title={Exemplo resolvido\ #1},left=6pt,right=6pt,top=3pt,bottom=3pt}
% Caixa de Atencao / erro comum
\newtcolorbox{atencao}[1][]{colback=WarnOrangeBg,colframe=WarnOrange,fonttitle=\bfseries,title={Atenção\ #1},left=6pt,right=6pt,top=3pt,bottom=3pt}
% Caixa "sai muitas vezes em exame"
\newtcolorbox{exame}[1][]{colback=ExamRedBg,colframe=ExamRed,fonttitle=\bfseries,title={Frequente em exame\ #1},left=6pt,right=6pt,top=3pt,bottom=3pt}

% Cada "Aula" (heading de nível 1, "#" em markdown -> \section) começa
% sempre numa página nova, para ficar fácil de encontrar a partir do índice.
% Sub-secções ("##", \subsection) NÃO forçam página nova.
\let\oldsection\section
\renewcommand{\section}{\clearpage\oldsection}

\pagestyle{fancy}
\fancyhf{}
\fancyhead[L]{\small\textcolor{gray}{\leftmark}}
\fancyhead[R]{\small\textcolor{gray}{\thepage}}
\renewcommand{\headrulewidth}{0.4pt}

\lstset{
  basicstyle=\ttfamily\small,
  breaklines=true,
  frame=single,
  columns=fullflexible,
  backgroundcolor=\color{gray!8}
}

\setlength{\parskip}{4pt}
\setlength{\parindent}{0pt}
